\section{Discussion, Limitations, and Ethical Considerations}
\label{sec:discussion}

\subsection{Architectural Synthesis and Insights}
The empirical findings across our 12 independent multi-seed runs provide three key insights for generalized facial forensics:
\begin{enumerate}[leftmargin=*]
    \item \textbf{Single-Frame LayerNorm PEFT is Remarkably Competitive:} Across both in-domain and out-of-domain evaluations, single-frame LayerNorm tuning (\texttt{smsdt\_s}, $103\text{K}$ parameters) achieves strong generalization ($0.8826$ on Celeb-DF v2 and $0.8841$ on DFD). Scaling trainable capacity to $8.5\text{M}$ (\texttt{smsdt\_st}) or $9.1\text{M}$ (\texttt{smsdt\_full}) yields mixed results ($+0.0056$ on Celeb-DF v2, $-0.0192$ on DFD), confirming that complex spatio-temporal heads do not reliably improve upon lightweight parameter-efficient tuning.
    \item \textbf{Optimization Recipe Dominates Architectural Additions:} The $+5.32\%$ Celeb-DF AUC gap between \texttt{smsdt\_s} and \texttt{gend\_ours} under identical parameter counts ($103{,}425$) demonstrates that loss balancing, gradient pacing, and curriculum sampling have a larger impact on generalization than adding millions of parameters to downstream heads.
    \item \textbf{FPR Inflation vs. AUC Trade-Off under Domain Shift:} While higher-capacity models can yield modest AUC improvements on compressed web benchmarks, they inflate the false alarm rate on pristine unseen faces (from $7.9\%$ to $22.5\%$). Standard summary metrics (AUC, HTER) conceal this deployment risk because benchmark evaluation splits are predominantly fake-skewed.
\end{enumerate}

\subsection{Limitations and Failure Modes}
Despite strong generalization, two clear boundary limitations remain:
\begin{itemize}[leftmargin=*]
    \item \textbf{Extreme Prior Asymmetries ($\pi \le 0.10$):} As audited in Table~\ref{tab:h3_prior}, when synthetic content prevalence drops to $10\%$, even a moderate false-positive rate produces a substantial volume of false alarms. While Saerens EM tracks the shift, extreme class imbalances require conservative operational thresholds to avoid degrading user trust.
    \item \textbf{Computational Footprint for Multi-Clip Inference:} Processing $K=8$ clips of $T=8$ frames requires $64$ total image passes per video. While suitable for asynchronous verification and forensic audit workflows, real-time edge streaming demands temporal pruning or progressive multi-clip cascading.
\end{itemize}

\subsection{Broader Impact and Ethical Considerations}
Facial forgery detection systems carry intrinsic dual-use implications. While critical for combating disinformation and protecting personal identity rights, automated detectors risk generating false accusations if deployed naively without calibration audits. Furthermore, forensic models can be inadvertently co-opted to construct more sophisticated adversarial generators via gradient-based perturbation. By open-sourcing our benchmark protocol, reproducible seeds, and complete statistical audit scripts, we encourage transparent, reproducible forensic research aligned with open scientific standards.

\section{Conclusion}
\label{sec:conclusion}
We presented a matched-pipeline empirical study investigating the generalization and operational stability of pretrained Vision Transformers for video deepfake detection. Across 12 independent multi-seed runs evaluated on FF++ c23, Celeb-DF-v2, and DFD, our findings demonstrate that: (1) single-frame LayerNorm PEFT achieves strong generalization without the computational overhead of spatio-temporal or multi-scale heads; (2) training recipe and loss formulation drive large performance deltas under identical parameter budgets; (3) operating point transfer exhibits minimal degradation ($g \le 0.0374 \ll 0.05$), falsifying catastrophic threshold collapse; and (4) higher-capacity heads introduce severe false positive inflation on unseen authentic faces, highlighting the necessity of false alarm budgeting in forensic evaluation.

